% SPDX-FileCopyrightText: 2026 Will Cook
% SPDX-License-Identifier: CC-BY-4.0
%
% One of the Erdős Problem Notes: a short standalone treatment of a single
% problem in the ErdosProblems expansion library.  Build with
% `tectonic erdos-269-three-prime-running-lcm.tex` or `make`.
\documentclass[11pt]{article}
% SPDX-FileCopyrightText: 2026 Will Cook
% SPDX-License-Identifier: CC-BY-4.0
%
% Shared preamble for the Erdős Problem Notes: the short, standalone,
% problem-owned manuscripts covering the expansion library `ErdosProblems`.
%
% One note owns one Erdős problem.  Each note repeats whatever context its own
% reader needs, so that a reader who arrives at a single problem never has to
% assemble it from the gateway paper.  Deliberate repetition across notes is the
% design, not an oversight.
%
% Source links resolve against one pinned formal-source commit, declared once
% below.  `scripts/check_problem_note_sources.py` verifies that every authored
% (file, line, declaration) triple names that exact declaration in the pinned
% snapshot, so a link cannot silently rot as later work moves lines.

\usepackage{paper-house-style}
\usepackage{array}
\usepackage{booktabs}
\usepackage{enumitem}
\setlist{topsep=0.35em,itemsep=0.2em}
\usepackage[colorlinks=true,linkcolor=linkinternal,urlcolor=linkexternal,
  citecolor=linkinternal]{hyperref}
\hypersetup{bookmarksnumbered=true,bookmarksopen=true,bookmarksopenlevel=1,
  pdfauthor={Will Cook},pdflang={en-GB}}

% ---------------------------------------------------------------------------
% Pinned formal source.  \commit names the checkpoint every link resolves
% against; the checker reads that snapshot rather than the working tree, so the
% printed coordinates stay correct however later waves move declarations.
% ---------------------------------------------------------------------------
\newcommand{\commit}{e5fd26a6d1b1a346430205eab5385b4c9565d004}
\newcommand{\commitshort}{e5fd26a6d1b1}
\newcommand{\repobase}{https://github.com/wcook04/plectis-lean-erdos249-257/blob/\commit}
\newcommand{\sourceurl}{https://github.com/wcook04/plectis-lean-erdos249-257/tree/\commit}
\newcommand{\PX}{ErdosProblems}

% Declaration names stay inside link targets.  The printed label is either a
% spaced, lower-case reading of the declaration name (\lref) or a short authored
% phrase (\lword); neither prints a module path or a raw Lean identifier.
\ExplSyntaxOn
\cs_new_protected:Npn \noteleanlabel #1
  {
    \tl_set:Nx \l_tmpa_tl { \tl_to_str:n {#1} }
    \regex_replace_all:nnN { _+ } { \c{space} } \l_tmpa_tl
    \regex_replace_all:nnN
      { ([a-z0-9])([A-Z]) }
      { \1\c{space}\2 }
      \l_tmpa_tl
    \tl_set:Nx \l_tmpa_tl { \tl_lower_case:n { \tl_use:N \l_tmpa_tl } }
    \tl_use:N \l_tmpa_tl
  }
\ExplSyntaxOff
\newcommand{\leanlabel}[1]{\noteleanlabel{#1}}
\newcommand{\lref}[3]{\href{\repobase/\PX/#1\#L#2}{\leanlabel{#3}}}
\newcommand{\lword}[4]{\href{\repobase/\PX/#1\#L#2}{#4}}
\newcommand{\lloc}[2]{\href{\repobase/\PX/#1\#L#2}{Lean source}}

% Links into a sibling library, named repository-relative.  The expansion
% library is implicit in \lref/\lword, because that is where problem-owned
% work lands.  Several problems have their headline theorem in the older
% reviewed #249/#257 corpus, and a note that could not name that library
% would have to paraphrase a checked statement instead of linking it.
\newcommand{\mref}[3]{\href{\repobase/#1\#L#2}{\leanlabel{#3}}}
\newcommand{\mword}[4]{\href{\repobase/#1\#L#2}{#4}}
\newcommand{\mloc}[2]{\href{\repobase/#1\#L#2}{Lean source}}
\emergencystretch=2em

\newtheorem{theorem}{Theorem}[section]
\newtheorem{proposition}[theorem]{Proposition}
\newtheorem{lemma}[theorem]{Lemma}
\newtheorem{corollary}[theorem]{Corollary}
\theoremstyle{definition}
\newtheorem{definition}[theorem]{Definition}
\newtheorem{example}[theorem]{Example}
\newtheorem{problem}[theorem]{Problem}
\theoremstyle{remark}
\newtheorem*{remark}{Remark}

\newcommand{\N}{\mathbb{N}}
\newcommand{\Npos}{\mathbb{N}_{>0}}
\newcommand{\Z}{\mathbb{Z}}
\newcommand{\Q}{\mathbb{Q}}
\newcommand{\R}{\mathbb{R}}
\newcommand{\lcm}{\operatorname{lcm}}
\newcommand{\ph}{\varphi}

% The series line printed under every note's title block.
\newcommand{\seriesline}{Erd\H{o}s Problem Notes}

% Production provenance.  The wording is fixed by the corpus provenance
% contract and reproduced verbatim in every manuscript in this repository, so
% the papers cannot drift into different accounts of how they were made.
\newcommand{\noteprovenance}{\provenancenote{\emph{Authorship and AI use.} The
prose, formal proofs, and software in this version were drafted and revised by
large-language-model agents under Will Cook's direction.  Cook set the
objectives and acceptance criteria, reviewed and authorised the manuscript, and
is responsible for its contents.  The agents are production tools, not authors.
See \emph{Statements and declarations}.}}

% The status paragraph every note carries, in its own words, before any result.
% Kernel checking establishes the exact Lean proposition; it does not establish
% that the proposition is the interesting one, that it is new, or that the
% Erdős problem is closed.
\newcommand{\statusboundary}{%
  \emph{Status.}  The problem treated here is open, and this note does not
  close it.  Every statement below marked as checked is a proposition that the
  pinned Lean kernel accepts from the sources this note links to, with no
  \texttt{sorry}, no added axiom, and no unchecked evaluation.  That is a claim
  about the formal statement, not about its mathematical interest, its
  novelty, or the original problem.  The unresolved obligations are named
  exactly, in their own section, and none of the finite computations,
  reductions, or no-go results here removes one of them.}

% Problem-specific source pin: #269's local-window consumer advances
% independently of the corpus default.
\renewcommand{\commit}{773d2a6ec5640b623c4a77262a82332d9ea85df0}
\renewcommand{\commitshort}{773d2a6ec564}
\hypersetup{
  pdftitle={The running least common multiple of a three-prime smooth prefix: Lean-checked structure for Erdős Problem 269},
  pdfsubject={Erdős Problem Notes: Problem 269},
  pdfkeywords={irrationality, least common multiple, smooth numbers, lattice sums, Lean 4},
}

% The three pure-power coordinates and the running height.
\newcommand{\hgt}{\operatorname{H}}
\newcommand{\lo}{\lambda}
\newcommand{\hi}{\eta}
\newcommand{\Lprefix}{\operatorname{L}}
\newcommand{\Kernel}{\operatorname{K}}

\runninghead{Erd\H{o}s \#269: three-prime running lcm}
\title{The Three-Prime Running Least Common Multiple}
\subtitle{Exact structure and the denominator-dependent local-window frontier}
\author{Will Cook}
\date{July 2026}

\begin{document}
\maketitle
\noteprovenance

\begin{abstract}
\noindent
Let $p,q,r$ be pairwise distinct primes and let $\Lprefix(x)$ denote the least
common multiple of all $\{p,q,r\}$-smooth numbers not exceeding $x$.  We prove
in Lean that
\[
 \Lprefix(x)=p^{\lfloor\log_p x\rfloor}q^{\lfloor\log_q x\rfloor}
 r^{\lfloor\log_r x\rfloor}
 \qquad(x\ge1),
\]
so the running least common multiple is the product of the three maximal pure
prime powers below the cutoff, and is at most $x^{3}$.  Consequently it is
constant on each cell where the three integer logarithms are constant; when
exactly one logarithm advances by one it is multiplied by exactly the
corresponding prime; and the first $n$ positive powers in the three channels,
together with the common origin $1$, give exactly $3n+1$ distinct jump values.
We prove an exact finite normal form in which a rectangular lattice sum of
reciprocal running values is grouped by its genuine height, with each
coefficient the cardinality of the corresponding fibre, and a uniform bound
$9\,\#\mathcal S\le(j+3)^{2}$ on the cardinality of a smooth exponent shell
whose sorted height coordinates sum to $j$.  At $\{2,3,5\}$ we exhibit the
smallest exact obstruction to separating the kernel as a product, the two sides
being $1/60$ and $1/12$.

Problem~\#269 is open and nothing here decides it.  We prove no irrationality
or transcendence statement in any three-prime case.  What is new at the
frontier is an exact consumer.  For a positive reduced denominator $B$ coprime
to $30$, a denominator-dependent cofinal local-window residue escape rules out
every positive integral carry obeying the matching bound.  Lean checks the
complete implication, including the zero-residue convention and all modulus
edge cases.  For the actual dyadic compression, Lean proves that each open
block $(2^a,2^{a+1})$ contains at most one pure $3$-power and at most one pure
$5$-power.  Hence its block radix belongs to the exact alphabet
$\{2,6,10,30\}$, rather than merely lying in a coarse bounded interval.  An
integer-only checker constructs the corresponding block digits, reproduces
three small certificates, and finds an escaping window of length at most $18$
for each of $106{,}666$ tested denominator/start pairs.  This is finite
experimental evidence, not a cofinal theorem.  What remains unproved is the
producer: bind the actual block digit and denominator-cleared carry to the
formal recurrence, then prove that these four-symbol windows escape cofinally.
\end{abstract}

\section{The problem}\label{sec:problem}

Let $P$ be a finite set of primes with $|P|\ge2$, and let
$a_1<a_2<\cdots$ enumerate the positive integers all of whose prime factors lie
in $P$.  Erd\H{o}s Problem~\#269 asks whether
\[
 \sum_{n\ge1}\frac{1}{[a_1,\ldots,a_n]}
\]
is irrational, where $[a_1,\ldots,a_n]$ is the least common
multiple~\cite[p.~65]{erdosgraham1980}\cite[p.~106]{erdos1988}.  Numbering and
status follow Bloom's catalogue~\cite{erdosproblems}.  The problem is open for
every finite $P$ with $|P|\ge2$.

Three things are known and fix the shape of the question.  The restriction
$|P|\ge2$ is necessary: for $P=\{p\}$ the enumeration is $a_n=p^{\,n-1}$, so
$[a_1,\ldots,a_n]=p^{\,n-1}$ and the sum is $p/(p-1)$, which is rational.  For
infinite $P$ the sum is always irrational, which Erd\H{o}s calls a simple
exercise~\cite[p.~106]{erdos1988}.  And in a letter of 1 January 1973 he
recorded that he could prove irrationality once duplicate summands are
removed~\cite{erdos1974letter}.

That last remark is the one this note is closest to.  The running least common
multiple is not injective in $n$: it is constant along stretches of the
enumeration and changes only at certain points.  Removing duplicate summands
means summing over the distinct values rather than over $n$.  Sections
\ref{sec:cells} and~\ref{sec:fibre} make that reindexing exact, in the case
$|P|=3$: they identify where the value is constant, by exactly what factor it
changes when it changes, and what multiplicity each distinct value carries.
They do not recover the irrationality proof Erd\H{o}s reported for the
de-duplicated sum, and they say nothing about the original sum.

Throughout, $p,q,r$ are pairwise distinct primes.  Call $n$ \emph{smooth} when
$n=p^{i}q^{j}r^{k}$ for some $i,j,k\ge0$; this is the
\lword{Erdos269/ThreePrimeRunningLcm.lean}{31}{smooth3Val}{smooth lattice
value}.  For $x\ge1$ write
\[
 \Lprefix(x)=\lcm\{\,n\le x:\ n\ \text{smooth}\,\},
 \qquad
 \hgt(x)=p^{\lfloor\log_p x\rfloor}\,q^{\lfloor\log_q x\rfloor}\,
 r^{\lfloor\log_r x\rfloor},
\]
the
\lword{Erdos269/ThreePrimeRunningLcm.lean}{53}{smoothPrefixLcm}{running least
common multiple} and the
\lword{Erdos269/ThreePrimeRunningLcm.lean}{36}{threePrimeHeight}{pure-power
height}, where $\lfloor\log_b x\rfloor$ is the integer logarithm, the largest
$e$ with $b^{e}\le x$.  Since $a_1,\ldots,a_n$ are exactly the smooth numbers
up to $a_n$, we have $[a_1,\ldots,a_n]=\Lprefix(a_n)$, and the summands of the
problem are the reciprocals of $\Lprefix$ along the enumeration.  The
reciprocal of the height at a smooth point is the
\lword{Erdos269/ThreePrimeRunningLcm.lean}{40}{threePrimeKernelQ}{lattice
kernel}
\[
 \Kernel(i,j,k)=\frac{1}{\hgt(p^{i}q^{j}r^{k})} .
\]

\paragraph{Scope of this note, and relation to prior work.}
We treat $|P|=3$ throughout, writing $P=\{p,q,r\}$, and the smallest instance
is $\{2,3,5\}$.  Nothing here is specific to three generators for its own sake;
three is where the coordinates stop separating, which
Theorem~\ref{res:rank} makes exact at the smallest instance.  The statement of
Problem~\#269 has been formalised before, as a conjecture with an unfilled
proof, in the \emph{Formal Conjectures} collection~\cite{formalconjectures269}.
That is a formal statement of the question.  The declarations described below
are propositions about the objects the question is posed over.  In relation to
the historical record, the result to foreground is not the elementary
running-lcm formula by itself but the checked denominator-dependent
local-window consumer: it isolates a precise producer whose proof would turn
the four-symbol radix structure into an irrationality contradiction.  No claim
of priority is made for any of these declarations.

\paragraph{Status of everything below.}
\statusboundary

\begin{table}[t]
\centering
\small
\begin{tabular}{@{}
  >{\raggedright\arraybackslash}p{0.34\linewidth}
  >{\raggedright\arraybackslash}p{0.18\linewidth}
  >{\raggedright\arraybackslash}p{0.40\linewidth}@{}}
\toprule
Statement & Status & Treatment here \\
\midrule
Irrationality in any three-prime case & Open & Not proved. \\
$\Lprefix(x)=\hgt(x)$ & Proved here & Theorem~\ref{res:lcm}; needs the primes distinct. \\
$\hgt(x)\le x^{3}$ & Proved here & Proposition~\ref{res:cube}. \\
Constancy on logarithmic cells & Proved here & Theorem~\ref{res:cell}. \\
Single-coordinate jump ratios & Proved here & Theorem~\ref{res:jump}. \\
Exactly $3n+1$ jump values & Proved here & Theorem~\ref{res:count}. \\
Height-fibre normal form & Proved here; finite & Theorem~\ref{res:fibre}. \\
Infinite limit of that expansion & Not proved & Section~\ref{sec:fibre}. \\
Shell multiplicity $9\,\#\mathcal S\le(j+3)^{2}$ & Proved here & Theorem~\ref{res:shell}. \\
Kernel is not separable at $\{2,3,5\}$ & Proved here & Theorem~\ref{res:rank}. \\
Actual dyadic block radix is in $\{2,6,10,30\}$ & Proved here & Theorem~\ref{res:dyadic-alphabet}. \\
Canonical least-positive-residue arithmetic & Proved here & Theorem~\ref{res:consumer}. \\
Denominator-dependent local-window consumer & Proved here & Theorem~\ref{res:windowconsumer}. \\
Dyadic-window scan for $B\le1000$ & Exact finite computation & $106{,}666$ pairs; no failures; not a theorem. \\
Local-window escape for the actual $\{2,3,5\}$ word & Open & Problem~\ref{prob:producer}. \\
\bottomrule
\end{tabular}
\end{table}

\paragraph{Reading map.}
Section~\ref{sec:lcm} identifies the running value.  Sections~\ref{sec:cells}
and~\ref{sec:fibre} develop its jump structure and the finite normal form,
Section~\ref{sec:shell} bounds fibre multiplicity, and Section~\ref{sec:rank}
records the obstruction to separating the kernel.  Section~\ref{sec:escape}
identifies the exact dyadic radix, exercises the finite consumer, and isolates
the remaining cofinal producer.  Every unproved hypothesis is marked as such.
Linked phrases open the corresponding Lean declaration at the pinned source
revision~\commitshort.

\medskip
\noindent\textbf{Keywords.} irrationality; least common multiple; smooth
numbers; lattice sums; Lean~4.
\qquad
\textbf{MSC 2020.} 11J72 (primary); 11A05, 11N25, 68V20 (secondary).

\section{The running value is a product of pure powers}\label{sec:lcm}

The smooth numbers up to $x$ are indexed by the exponent triples $(i,j,k)$ with
$i\le\lfloor\log_p x\rfloor$, $j\le\lfloor\log_q x\rfloor$,
$k\le\lfloor\log_r x\rfloor$ and $p^{i}q^{j}r^{k}\le x$: the coordinate bounds
make the index set finite, and the last condition is the actual constraint.
This is the
\lword{Erdos269/ThreePrimeRunningLcm.lean}{46}{smoothPrefixExponents}{smooth
prefix index set}.  Both conditions matter: the coordinate box is strictly
larger than the prefix, since a product of three large pure powers can exceed
$x$ while each factor does not.

\begin{theorem}[the running least common multiple]\label{res:lcm}
Let $p,q,r$ be pairwise distinct primes and $x\ge1$.  Then
$\Lprefix(x)=\hgt(x)$.
\end{theorem}

\begin{proof}
For divisibility in one direction, every smooth $n\le x$ has exponents bounded
by the corresponding integer logarithms, so $n\mid\hgt(x)$, and hence
$\Lprefix(x)\mid\hgt(x)$.  For the other, the three pure powers
$p^{\lfloor\log_p x\rfloor}$, $q^{\lfloor\log_q x\rfloor}$ and
$r^{\lfloor\log_r x\rfloor}$ are themselves smooth numbers not exceeding $x$,
so each divides $\Lprefix(x)$.  Distinct primes have coprime powers, so their
product divides $\Lprefix(x)$ as well.  The two divisibilities give equality.
\end{proof}

\noindent Formalised as the
\lword{Erdos269/ThreePrimeRunningLcm.lean}{123}{smoothPrefixLcm_eq_threePrimeHeight}{running-lcm
identity}, from the
\lword{Erdos269/ThreePrimeRunningLcm.lean}{77}{smoothPrefixLcm_dvd_threePrimeHeight}{divisibility
into the height} and the three membership statements for the pure powers,
namely the
\lword{Erdos269/ThreePrimeRunningLcm.lean}{84}{pureFirst_mem_smoothPrefixExponents}{first},
\lword{Erdos269/ThreePrimeRunningLcm.lean}{97}{pureSecond_mem_smoothPrefixExponents}{second},
and
\lword{Erdos269/ThreePrimeRunningLcm.lean}{110}{pureThird_mem_smoothPrefixExponents}{third}
pure-power memberships.

The hypothesis that the primes are pairwise distinct is used exactly once, in
the coprimality step, and it is not decorative: without it the three pure
powers need not have coprime orders and their product need not divide the least
common multiple.  The identity is what makes every later statement about
$\Lprefix$ computable from three integer logarithms.

\begin{proposition}\label{res:cube}
For every $x\ge1$, $\hgt(x)\le x^{3}$.
\end{proposition}

\begin{proof}
Each of the three factors is a power of its base not exceeding $x$.
\end{proof}

\noindent Formalised as the
\lword{Erdos269/ThreePrimeRunningLcm.lean}{430}{threePrimeHeight_le_cube}{cubic
majorant}, which needs no primality.  The exponent $3$ is the number of
generating primes, and the bound is the reason the reciprocal kernel is
comparable to $x^{-3}$ rather than to $x^{-1}$.

\section{Cells, jumps, and their count}\label{sec:cells}

Say that $x$ and $y$ lie in the same \emph{logarithmic cell} when
$\lfloor\log_b x\rfloor=\lfloor\log_b y\rfloor$ for each of $b=p,q,r$: the
\lword{Erdos269/ThreePrimeRunningLcm.lean}{163}{SameThreePrimeLogCell}{cell
relation}.

\begin{theorem}[constancy on cells]\label{res:cell}
If $x,y\ge1$ lie in the same logarithmic cell, then
$\Lprefix(x)=\Lprefix(y)$.  The same holds for the kernel at two smooth points
whose values lie in one cell.
\end{theorem}

\begin{proof}
Immediate from Theorem~\ref{res:lcm}, since $\hgt$ depends on $x$ only
through the three integer logarithms.
\end{proof}

\noindent Formalised as the
\lword{Erdos269/ThreePrimeRunningLcm.lean}{179}{smoothPrefixLcm_eq_of_sameLogCell}{cell
constancy of the running value}, the
\lword{Erdos269/ThreePrimeRunningLcm.lean}{169}{threePrimeHeight_eq_of_sameLogCell}{cell
constancy of the height}, and the
\lword{Erdos269/ThreePrimeRunningLcm.lean}{191}{threePrimeKernelQ_eq_of_sameLogCell}{cell
constancy of the kernel}.  The running value therefore changes only when one of
the three logarithms changes, and the next theorem says by exactly how much.

\begin{theorem}[single-coordinate jump ratios]\label{res:jump}
Let $x,y\ge1$.  If $\lfloor\log_p y\rfloor=\lfloor\log_p x\rfloor+1$ while the
other two logarithms agree, then $\Lprefix(y)=p\,\Lprefix(x)$; and similarly
with $q$ or $r$ in place of $p$.
\end{theorem}

\begin{proof}
By Theorem~\ref{res:lcm} both sides are heights, and one factor of the height
gains one in its exponent while the others are unchanged.
\end{proof}

\noindent Formalised as the
\lword{Erdos269/ThreePrimeRunningLcm.lean}{326}{smoothPrefixLcm_firstLogStep}{first},
\lword{Erdos269/ThreePrimeRunningLcm.lean}{339}{smoothPrefixLcm_secondLogStep}{second},
and
\lword{Erdos269/ThreePrimeRunningLcm.lean}{352}{smoothPrefixLcm_thirdLogStep}{third}
coordinate steps, over the corresponding statements for the height alone, which
need no primality: the
\lword{Erdos269/ThreePrimeRunningLcm.lean}{294}{threePrimeHeight_firstLogStep}{first
height step}, the
\lword{Erdos269/ThreePrimeRunningLcm.lean}{305}{threePrimeHeight_secondLogStep}{second},
and the
\lword{Erdos269/ThreePrimeRunningLcm.lean}{316}{threePrimeHeight_thirdLogStep}{third}.

The jump values are therefore the pure prime powers, and they do not collide
across channels.

\begin{theorem}[jump count]\label{res:count}
Let $n\ge0$.  The set of the first $n$ positive powers of $p$, of $q$, and of
$r$ has exactly $3n$ elements, and adjoining the common origin $1$ gives
exactly $3n+1$.
\end{theorem}

\begin{proof}
Within one channel the powers $b,b^{2},\ldots,b^{n}$ are distinct because
$b\ge2$.  Across two channels a common value would be a positive power of two
distinct primes, which unique factorisation forbids.  Finally $1$ is not a
positive power of any prime.
\end{proof}

\noindent Formalised as the
\lword{Erdos269/ThreePrimeRunningLcm.lean}{249}{threePrimePositiveJumpSet_card}{positive
jump count} and the
\lword{Erdos269/ThreePrimeRunningLcm.lean}{278}{threePrimeJumpSetWithOrigin_card}{jump
count with the origin}, over the
\lword{Erdos269/ThreePrimeRunningLcm.lean}{208}{positivePrimePowers_card}{channel
cardinality}, the
\lword{Erdos269/ThreePrimeRunningLcm.lean}{229}{positivePrimePowers_disjoint}{channel
disjointness}, and the
\lword{Erdos269/ThreePrimeRunningLcm.lean}{219}{one_not_mem_positivePrimePowers}{exclusion
of the origin}; the channels themselves are the
\lword{Erdos269/ThreePrimeRunningLcm.lean}{204}{positivePrimePowers}{positive
power sets}.  The exponent $0$ is omitted from each channel because it is the
shared initial value, which is why the origin is counted once rather than three
times.

\section{The height-fibre normal form}\label{sec:fibre}

Grouping a lattice sum by the value of the running least common multiple turns
it into a sum over heights whose coefficients are multiplicities.  We prove
this exactly, on a finite rectangular box.

Write $B(h_p,h_q,h_r)$ for the box of exponent triples with $i\le h_p$,
$j\le h_q$, $k\le h_r$, the
\lword{Erdos269/ThreePrimeRunningLcm.lean}{368}{smoothExponentBox}{exponent
box}, and for a height $H$ write $F(H)$ for the set of points of the box whose
running height equals $H$, the
\lword{Erdos269/ThreePrimeRunningLcm.lean}{377}{smoothHeightFiber}{height
fibre}.

\begin{theorem}[finite normal form]\label{res:fibre}
For every box,
\[
 \sum_{(i,j,k)\in B}\Kernel(i,j,k)
 =\sum_{H}\frac{\#F(H)}{H},
\]
the outer sum ranging over the heights actually attained on $B$.
\end{theorem}

\begin{proof}
Partition $B$ into the fibres of the height map.  On $F(H)$ every summand is
$1/H$ by definition of the kernel, so the fibre contributes $\#F(H)/H$.
\end{proof}

\noindent Formalised as the
\lword{Erdos269/ThreePrimeRunningLcm.lean}{406}{finiteSmoothKernelSum_groupedByHeight}{height-fibre
normal form}, over the
\lword{Erdos269/ThreePrimeRunningLcm.lean}{384}{smoothHeightFiber_kernel_sum}{fibre
sum}, with the height of a lattice point given by the
\lword{Erdos269/ThreePrimeRunningLcm.lean}{373}{smoothPointHeight}{point
height}.

Together with Theorem~\ref{res:cell} this is the finite core of the ordered
prime-power jump expansion: the value is constant on cells, the cells are
indexed by heights, and the coefficient of a height is the number of lattice
points it collects.  The passage to the infinite sum, and the explicit ordering
of the pure powers that would make the expansion a series in the jumps, are not
proved here.  Theorem~\ref{res:fibre} is an identity between two finite sums,
and it is stated over the full box rather than the smooth prefix, so it is not
a statement about $\Lprefix$ at a cutoff.

The tail dynamics that such an expansion would feed is recorded in the same
module as a one-step map, $\tau(b,d,s)=b(s-d)$, the
\lword{Erdos269/ThreePrimeRunningLcm.lean}{487}{tailStateStep}{variable-base
tail step}, together with the rewriting
\lword{Erdos269/ThreePrimeRunningLcm.lean}{490}{tailStateStep_eq}{that names its
expanded form}.  No orbit of this map is analysed.

\section{Shell multiplicity}\label{sec:shell}

A tail estimate needs to know how many lattice points can share a short
multiplicative interval.  Write $\mathcal S$ for the set of exponent triples in
a box $B(h_p,h_q,h_r)$ whose smooth value lies in $[\lo,\hi)$, the
\lword{Erdos269/ThreePrimeRunningLcm.lean}{500}{smoothExponentShell}{smooth
exponent shell}.

\begin{lemma}[uniqueness in a short interval]\label{res:short}
Let $b\ge1$ and suppose $\hi\le b\,\lo$.  If $b^{a}w$ and $b^{a'}w$ both lie in
$[\lo,\hi)$, then $a=a'$.
\end{lemma}

\begin{proof}
If $a<a'$ then $\hi\le b\,\lo\le b\cdot b^{a}w=b^{a+1}w\le b^{a'}w<\hi$, which
is impossible; the case $a>a'$ is symmetric.
\end{proof}

\noindent Formalised as the
\lword{Erdos269/ThreePrimeRunningLcm.lean}{509}{exponent_unique_in_short_interval}{short-interval
uniqueness}.  The hypothesis is that the interval has multiplicative width at
most $b$; then one coordinate is determined by the other two, and projecting it
away is injective.

\begin{proposition}\label{res:drop}
If $\hi\le r\,\lo$ then $\#\mathcal S\le(h_p+1)(h_q+1)$; if $\hi\le p\,\lo$ then
$\#\mathcal S\le(h_q+1)(h_r+1)$.
\end{proposition}

\noindent Formalised as the
\lword{Erdos269/ThreePrimeRunningLcm.lean}{585}{smoothExponentShell_card_le_dropThird}{third-coordinate
projection} and the
\lword{Erdos269/ThreePrimeRunningLcm.lean}{543}{smoothExponentShell_card_le_dropFirst}{first-coordinate
projection}.

\begin{theorem}[quadratic multiplicity bound]\label{res:shell}
Suppose $\hi\le r\,\lo$ and $h_p\le h_q\le h_r$ with $h_p+h_q+h_r=j$.  Then
\[
 9\,\#\mathcal S\le(j+3)^{2}.
\]
\end{theorem}

\begin{proof}
By Proposition~\ref{res:drop}, $\#\mathcal S\le(h_p+1)(h_q+1)$.  It therefore
suffices to prove that $a\le b\le c$ with $a+b+c=j$ gives
$9(a+1)(b+1)\le(j+3)^{2}$.  From $a\le b\le c$ we get $a+2b\le j$, so it is
enough that $9(a+1)(b+1)\le(a+2b+3)^{2}$; writing $b=a+d$ with $d\ge0$, the
difference of the two sides is $d(3a+4d+3)\ge0$.
\end{proof}

\noindent Formalised as the
\lword{Erdos269/ThreePrimeRunningLcm.lean}{646}{smoothExponentShell_card_quadratic}{quadratic
shell bound}, over the
\lword{Erdos269/ThreePrimeRunningLcm.lean}{629}{sorted_pair_quadratic}{sorted
quadratic estimate}.  The constant $9$ is the square of the number of
generating primes and appears because the bound is the arithmetic--geometric
comparison for a sum of three sorted coordinates; sorting is a hypothesis, not
a normalisation, since the shell itself is not symmetric in the three bases.

The bound is uniform in $\lo$ and $\hi$ subject to the width condition, and it
is stated for the actual filtered shell rather than for a lattice model of it.
It is an input to a tail estimate and is not itself one: no series is bounded
here.

\section{The kernel does not separate}\label{sec:rank}

A recurring hope is to write the three-prime kernel as a product of functions
of the separate coordinates, which would reduce the problem to one-dimensional
criteria.  At the smallest instance this fails, and the failure is exact.

\begin{theorem}[non-separability at $\{2,3,5\}$]\label{res:rank}
With $(p,q,r)=(2,3,5)$,
\[
 \Kernel(0,0,0)\,\Kernel(1,1,0)=\tfrac1{60}
 \ne\tfrac1{12}=\Kernel(1,0,0)\,\Kernel(0,1,0).
\]
\end{theorem}

\begin{proof}
The four values are $\Kernel(0,0,0)=1$, $\Kernel(1,0,0)=1/2$,
$\Kernel(0,1,0)=1/6$ and $\Kernel(1,1,0)=1/60$, computed from
$\hgt(1)=1$, $\hgt(2)=2$, $\hgt(3)=2\cdot3=6$ and
$\hgt(6)=4\cdot3\cdot5=60$.
\end{proof}

\noindent Formalised as the
\lword{Erdos269/ThreePrimeRunningLcm.lean}{479}{kernel_235_not_rankOne}{non-separation
witness}, over the four exact values, the
\lword{Erdos269/ThreePrimeRunningLcm.lean}{443}{kernel_235_origin}{origin},
\lword{Erdos269/ThreePrimeRunningLcm.lean}{450}{kernel_235_two}{value at two},
\lword{Erdos269/ThreePrimeRunningLcm.lean}{458}{kernel_235_three}{value at
three}, and
\lword{Erdos269/ThreePrimeRunningLcm.lean}{467}{kernel_235_six}{value at six}.
The height at $6$ is $60$ rather than $6$ because the maximal pure powers below
$6$ are $4$, $3$ and $5$: the running least common multiple at a smooth cutoff
sees powers of the other primes that the cutoff itself does not contain.  That
is the mechanism behind the non-separation.

A nonzero two-by-two minor rules out writing the kernel as
$f(i)g(j)h(k)$ on this box, and hence rules out any argument that first
separates the coordinates and then treats them one at a time.  It does not
bound the rank of the kernel from below beyond one, and it is not an
independence or irrationality statement.

\section{The exact local-window frontier}\label{sec:escape}

\subsection{The actual four-symbol dyadic radix}

Compress the jump word between consecutive powers of two: a block starts just
after $2^a$, includes every pure $3$- or $5$-power strictly between $2^a$ and
$2^{a+1}$, and ends with the jump at $2^{a+1}$.  A channel cannot occur twice
inside one block.  Indeed, if
\[
  2^a<p^e,p^f<2^{a+1}\qquad(p\ge2),
\]
then the ratio between the interval endpoints is $2\le p$, so strict
monotonicity of the powers forces $e=f$.  This is the
\lword{Erdos269/ThreePrimeRunningLcm.lean}{668}{dyadicInternalPower_exponent_unique}{checked
internal-power uniqueness lemma}.

Let $\beta_a$ be the product of the terminal dyadic factor $2$, a factor $3$
when the block contains an internal $3$-power, and a factor $5$ when it
contains an internal $5$-power.

\begin{theorem}[the dyadic block alphabet]\label{res:dyadic-alphabet}
For every $a$,
\[
  \beta_a\in\{2,6,10,30\};
  \qquad\text{in particular}\qquad 2\le\beta_a\le30.
\]
\end{theorem}

\begin{proof}
Internal-power uniqueness leaves two independent yes/no choices, one for the
$3$-channel and one for the $5$-channel.  Multiplying the terminal factor $2$
by the selected channel factors gives exactly the displayed four cases.
\end{proof}

\noindent The definition is the
\lword{Erdos269/ThreePrimeRunningLcm.lean}{688}{dyadicBlockBase235}{dyadic
block base}; Lean checks both the
\lword{Erdos269/ThreePrimeRunningLcm.lean}{699}{dyadicBlockBase235_cases}{exact
four-case alphabet} and the
\lword{Erdos269/ThreePrimeRunningLcm.lean}{711}{dyadicBlockBase235_mem_interval}{bounded-radix
consequence}.  Thus radix growth is no longer an open part of the producer.
What is not yet checked in Lean is the identification of the corresponding
block digit with the multiplicity forcing in the original series.

\paragraph{Exact finite consumer experiment.}
The integer-only
\href{\sourceurl/scripts/check_erdos269_dyadic_windows.py}{dyadic-window
checker} constructs the ordered pure-power jumps, the block bases, and the
block digits from exact multiplicity counts.  It reproduces the following
certificates; the columns are denominator $B$, dyadic start $a$, window length
$h$, endpoint jump index $n$, window base $W$, forcing $F$, least positive
residue $R=\operatorname{lpr}_{W}(-BF)$, and short bound $K$.
\[
\begin{array}{c|c|c|c|r|r|r|r}
B&a&h&n&W&F&R&K\\ \hline
1&1&2&4&60&47&13&9\\
7&1&3&6&360&289&137&95\\
16&1&4&9&10800&8735&640&352
\end{array}
\]
A fresh scan over every $B\le1000$ coprime to $30$ and every
$100\le a\le500$ tested $106{,}666$ pairs.  In every case a window of length
at most $18$ made both $W>K$ and
$\operatorname{lpr}_{W}(-BF)>K$; the largest first successful length was
$14$.  The computation uses integers only and is reproducible from the pinned
checker.  Its claim ceiling is deliberately finite: neither the scan nor the
three displayed certificates proves escape for unbounded $B$ or cofinally
many starts.

The residue route can now be stated without hiding either the denominator or
the consumer.  For an integer radix word $b_n$ and forcing word $m_n$, define
\[
\begin{aligned}
 W_{\ell,0}&=1,
&W_{\ell,h+1}&=b_{\ell+h}W_{\ell,h},\\
 F_{\ell,0}&=0,
&F_{\ell,h+1}&=b_{\ell+h}F_{\ell,h}+m_{\ell+h}.
\end{aligned}
\]
These are the
\lword{Erdos269/RestrictedFloorSum.lean}{16}{windowBase}{window base} and
\lword{Erdos269/RestrictedFloorSum.lean}{21}{windowForcing}{window forcing}.
If an integral carry satisfies
\[
 d_{n+1}=b_n d_n-Bm_n ,
\]
then induction gives the exact division-free identity
\[
 d_{\ell+h}=W_{\ell,h}d_\ell-BF_{\ell,h};
\]
this is the
\lword{Erdos269/RestrictedFloorSum.lean}{52}{integralCarry_window}{checked
window identity}.

For $C>0$, let $\operatorname{lpr}_C(x)\in\{1,\ldots,C\}$ be the least
positive representative of $x\bmod C$, with a zero residue represented by
$C$.  This is the
\lword{Erdos269/ResidueEscape.lean}{19}{leastPositiveResidue}{canonical
representative}.  Lean checks both its
\lword{Erdos269/ResidueEscape.lean}{24}{leastPositiveResidue_pos_le}{positive
range} and its
\lword{Erdos269/ResidueEscape.lean}{45}{leastPositiveResidue_modEq}{congruence
to the source integer}.  This convention matters: replacing
$\operatorname{lpr}_C(x)$ by $|x|$ would not be a modular statement.

Let $K(B,n)$ be the bound on the reduced carry at denominator $B$.  Define
\emph{cofinal local-window escape} to mean that for every $B>0$ coprime to
$30$ and every $\ell_0$, there are $\ell\ge\ell_0$ and $h>0$ such that
\[
 C_{\ell,h}:=|W_{\ell,h}|>0
 \quad\text{and}\quad
 K(B,\ell+h)<
 \operatorname{lpr}_{C_{\ell,h}}(-BF_{\ell,h}).
\tag{E}\label{eq:escape}
\]
The quantifier over $B$ is part of the proposition, and so is the dependence
of $K$ on $B$; suppressing it gives the wrong expert interface.  The exact
unproved proposition is the
\lword{Erdos269/RestrictedFloorSum.lean}{69}{CofinalLocalWindowEscape}{cofinal
local-window producer}.

\begin{theorem}[the finite residue contradiction]\label{res:consumer}
Let $C>0$, $c>0$, and $|c|\le K$.  If
$c\equiv x\pmod C$ and $K<\operatorname{lpr}_C(x)$, then the hypotheses are
contradictory.
\end{theorem}

\begin{proof}
The canonical representative lies in $\{1,\ldots,C\}$.  The inequalities put
$c$ strictly between $0$ and $C$.  If $x\equiv0\pmod C$, its positive
representative is $C$ while $c\bmod C=c\ne0$.  Otherwise both $c$ and
$\operatorname{lpr}_C(x)$ are their own residues and congruence makes them
equal, contradicting $|c|\le K<\operatorname{lpr}_C(x)$.
\end{proof}

\noindent The natural-state version is the
\lword{Erdos269/ResidueEscape.lean}{71}{no_bounded_positive_state_of_residue_escape}{finite
consumer}; the integer carry version is the
\lword{Erdos269/ResidueEscape.lean}{104}{no_bounded_positive_int_state_of_leastPositiveResidue}{least-positive-residue
consumer}.

\begin{theorem}[the checked local-window consumer]\label{res:windowconsumer}
Assume the cofinal escape property~\eqref{eq:escape}.  Fix $B>0$ coprime to
$30$.  There is no integral sequence $d_n$ satisfying simultaneously
\[
 d_{n+1}=b_nd_n-Bm_n,\qquad d_n>0,\qquad
 |d_n|\le K(B,n)\quad(n\ge0).
\]
\end{theorem}

\begin{proof}
Choose one window supplied by~\eqref{eq:escape}.  The checked window identity
gives
\[
 d_{\ell+h}\equiv -BF_{\ell,h}\pmod{|W_{\ell,h}|}.
\]
The endpoint state is positive and at most $K(B,\ell+h)$, whereas the
canonical positive residue of the right-hand side is larger than this bound.
Theorem~\ref{res:consumer} is the contradiction.
\end{proof}

\noindent This is formalised as the
\lword{Erdos269/RestrictedFloorSum.lean}{83}{no_positive_reducedCarry_of_cofinalLocalWindowEscape}{reduced-carry
extinction theorem}.  Coprimality with $30$ is used by the open producer to
select a window; once a window has been emitted, the finite contradiction
does not use it.  The theorem also treats the edge cases explicitly:
$|W_{\ell,h}|=0$ is excluded, a zero residue is represented by the full
modulus, and positivity prevents the endpoint carry from being zero.

\subsection{Candidate mechanisms and disconfirmers}

The most direct mechanism is now four-symbol anti-concentration.  The radix
word is constrained to $\{2,6,10,30\}$, so the surviving arithmetic lies in
the block digit and its carry into the residue class.  One seeks a cofinal
family of blocks for which the modulus $|W_{\ell,h}|$ grows faster than the
carry window and the weighted forcing $-BF_{\ell,h}$ does not remain trapped
in the two short arcs adjacent to $0$ modulo that modulus.  A theorem for one
cofinal family is enough; no equidistribution of all windows is required.

Four shortcuts are already disconfirmed.  Merely proving
$2\le b_a\le30$ loses the exact channel alphabet.  A fixed finite list of
denominator certificates does not address the universal $B$-quantifier.  A
bound independent of $B$ is not the bound delivered by denominator clearing.
Finally, large absolute forcing is irrelevant unless its \emph{least positive
residue} is large.  These are not stylistic cautions: each removes a false
version of~\eqref{eq:escape}.

\section{Specialist handoff and surviving questions}\label{sec:open}

\begin{problem}[the exact surviving producer]\label{prob:producer}
For the actual ordered $\{2,3,5\}$ jump expansion, identify the formal radix
word with $\beta_a\in\{2,6,10,30\}$, define and verify the corresponding block
digit $m_a$, derive the denominator-cleared bound $K(B,a)$ from a hypothetical
rational value, and prove~\eqref{eq:escape} for every $B>0$ coprime to $30$
and every starting index.  Equivalently, prove a cofinal subfamily of windows
on which
\[
 \operatorname{lpr}_{|W_{\ell,h}|}(-BF_{\ell,h})>K(B,\ell+h).
\]
\end{problem}

\paragraph{Wanted.}
A theorem proving Problem~\ref{prob:producer}, or a strictly weaker
four-symbol block statement that still emits one escaping window beyond every
level for each fixed $B$.  The answer must state the block digit, all
dependence on $B$, and the endpoint at which the carry bound is evaluated.

\paragraph{Already checked.}
Lean checks the exact dyadic radix alphabet, the window recurrence, the
least-positive-residue interval and congruence lemmas, the
denominator-dependent producer type, and the complete
producer-to-contradiction theorem.  The structural height, jump, fibre,
quadratic shell, and non-separation results earlier in this note are also
checked.  The exact checker constructs actual block digits and supplies the
finite evidence reported above.  No Lean theorem yet identifies those digits
with the formal forcing word, derives $K(B,a)$ from a rational value, or
proves~\eqref{eq:escape}.

\paragraph{Most useful answer.}
First close the exact bridge from the original summands to the block digit
used by the checker and to the denominator-cleared bound.  Then give an
explicit cofinal rule in the alphabet $\{2,6,10,30\}$, together with a lower
bound for the canonical residue and an upper bound for $K(B,\ell+h)$.  A
negative answer is also valuable if it constructs infinitely many
$(B,\ell_0)$ for which every proposed window family fails; that would retire
the present producer rather than merely weaken it.

\paragraph{If answered.}
Theorem~\ref{res:windowconsumer} immediately rules out the corresponding
positive reduced carry.  After the separate rational-value-to-actual-carry
instantiation is checked, this supplies the missing contradiction for the
three-prime case.  Until both pieces are present, Problem~\#269 remains open.

\begin{problem}[direct nonintegral tails]
As an alternative to local windows, prove directly that for every positive
$B$ coprime to $30$ the denominator-cleared actual dyadic tail is nonintegral
cofinally.  Such a theorem would bypass~\eqref{eq:escape}, but it must still
feed the same rational-carry consumer.
\end{problem}

\begin{problem}[a genuinely three-dimensional analytic theorem]
Prove a target-faithful analytic statement for the phase cocycle attached to
the three logarithmic coordinates.  By Theorem~\ref{res:rank}, a preliminary
separation into one-dimensional factors is unavailable even at $\{2,3,5\}$.
\end{problem}

\pdfbookmark[1]{Statements and declarations}{statements-declarations}
\subsection*{Statements and declarations}

\paragraph{Artefact and data availability.} The
\href{\sourceurl}{pinned formal-source revision} contains the Lean sources, the
fixed toolchain, the library manifest, and the exact dyadic-window checker used
in the finite experiment.  This manuscript provides navigation rather than
proof authority.

\paragraph{Declaration of generative AI use.} Large-language-model agents were
used throughout development to draft and revise prose, formal proofs, and
software.  The author set the objectives and acceptance criteria, selected and
reviewed the claims, and approved the published version.  The author assumes
responsibility for the accuracy, interpretation, and presentation of the work.
Generative systems are production tools; they are not authors and supply no
independent authority.  Lean checks each proof term against the fixed library
version, and the sources linked here contain no proof placeholders and no
project-defined axioms; Lean does not authorise the exposition, the citation
choices, or the interpretation, for which the author remains responsible.

\paragraph{Funding and competing interests.} This work received no external
funding.  The author declares no competing interests.

\paragraph{Acknowledgements.} The problem numbering and status follow the
Erd\H{o}s Problems catalogue maintained by Thomas Bloom~\cite{erdosproblems}.

\appendix
\section{Guide to the formal sources}\label{app:index}

Each linked phrase opens its Lean declaration at the pinned source
revision~\commitshort.  The running-LCM structure, residue arithmetic, and
local-window bridge occupy three separate modules.  Three distinctions are
worth carrying into the source.  The height statements hold for arbitrary
bases, while the statements about $\Lprefix$ need the three primes to be
distinct.  The normal form of Section~\ref{sec:fibre} is a finite identity over
a rectangular box, not a convergence theorem.  And the formal cofinal-escape
predicate is an unproved proposition consumed by a theorem; its application
to the actual $\{2,3,5\}$ word is not asserted.

% The exhaustive source manifest is retained in the authored TeX for automated
% validation, and kept outside the printed note.
\iffalse

\begin{tabular}{@{}ll@{}}
informal statement & linked Lean source\\
smooth lattice value & \lrefx{Erdos269/ThreePrimeRunningLcm.lean}{31}{smooth3Val}\\
pure-power height & \lrefx{Erdos269/ThreePrimeRunningLcm.lean}{36}{threePrimeHeight}\\
lattice kernel & \lrefx{Erdos269/ThreePrimeRunningLcm.lean}{40}{threePrimeKernelQ}\\
smooth prefix index set & \lrefx{Erdos269/ThreePrimeRunningLcm.lean}{46}{smoothPrefixExponents}\\
running least common multiple & \lrefx{Erdos269/ThreePrimeRunningLcm.lean}{53}{smoothPrefixLcm}\\
prefix value divides the height & \lrefx{Erdos269/ThreePrimeRunningLcm.lean}{59}{smooth3Val_dvd_threePrimeHeight_of_mem}\\
divisibility into the height & \lrefx{Erdos269/ThreePrimeRunningLcm.lean}{77}{smoothPrefixLcm_dvd_threePrimeHeight}\\
first pure-power membership & \lrefx{Erdos269/ThreePrimeRunningLcm.lean}{84}{pureFirst_mem_smoothPrefixExponents}\\
second pure-power membership & \lrefx{Erdos269/ThreePrimeRunningLcm.lean}{97}{pureSecond_mem_smoothPrefixExponents}\\
third pure-power membership & \lrefx{Erdos269/ThreePrimeRunningLcm.lean}{110}{pureThird_mem_smoothPrefixExponents}\\
running-lcm identity & \lrefx{Erdos269/ThreePrimeRunningLcm.lean}{123}{smoothPrefixLcm_eq_threePrimeHeight}\\
cell relation & \lrefx{Erdos269/ThreePrimeRunningLcm.lean}{163}{SameThreePrimeLogCell}\\
cell constancy of the height & \lrefx{Erdos269/ThreePrimeRunningLcm.lean}{169}{threePrimeHeight_eq_of_sameLogCell}\\
cell constancy of the running value & \lrefx{Erdos269/ThreePrimeRunningLcm.lean}{179}{smoothPrefixLcm_eq_of_sameLogCell}\\
cell constancy of the kernel & \lrefx{Erdos269/ThreePrimeRunningLcm.lean}{191}{threePrimeKernelQ_eq_of_sameLogCell}\\
positive power sets & \lrefx{Erdos269/ThreePrimeRunningLcm.lean}{204}{positivePrimePowers}\\
channel cardinality & \lrefx{Erdos269/ThreePrimeRunningLcm.lean}{208}{positivePrimePowers_card}\\
exclusion of the origin & \lrefx{Erdos269/ThreePrimeRunningLcm.lean}{219}{one_not_mem_positivePrimePowers}\\
channel disjointness & \lrefx{Erdos269/ThreePrimeRunningLcm.lean}{229}{positivePrimePowers_disjoint}\\
positive jump channels & \lrefx{Erdos269/ThreePrimeRunningLcm.lean}{243}{threePrimePositiveJumpSet}\\
positive jump count & \lrefx{Erdos269/ThreePrimeRunningLcm.lean}{249}{threePrimePositiveJumpSet_card}\\
jump set with the origin & \lrefx{Erdos269/ThreePrimeRunningLcm.lean}{273}{threePrimeJumpSetWithOrigin}\\
jump count with the origin & \lrefx{Erdos269/ThreePrimeRunningLcm.lean}{278}{threePrimeJumpSetWithOrigin_card}\\
first height step & \lrefx{Erdos269/ThreePrimeRunningLcm.lean}{294}{threePrimeHeight_firstLogStep}\\
second height step & \lrefx{Erdos269/ThreePrimeRunningLcm.lean}{305}{threePrimeHeight_secondLogStep}\\
third height step & \lrefx{Erdos269/ThreePrimeRunningLcm.lean}{316}{threePrimeHeight_thirdLogStep}\\
first coordinate step & \lrefx{Erdos269/ThreePrimeRunningLcm.lean}{326}{smoothPrefixLcm_firstLogStep}\\
second coordinate step & \lrefx{Erdos269/ThreePrimeRunningLcm.lean}{339}{smoothPrefixLcm_secondLogStep}\\
third coordinate step & \lrefx{Erdos269/ThreePrimeRunningLcm.lean}{352}{smoothPrefixLcm_thirdLogStep}\\
exponent box & \lrefx{Erdos269/ThreePrimeRunningLcm.lean}{368}{smoothExponentBox}\\
point height & \lrefx{Erdos269/ThreePrimeRunningLcm.lean}{373}{smoothPointHeight}\\
height fibre & \lrefx{Erdos269/ThreePrimeRunningLcm.lean}{377}{smoothHeightFiber}\\
fibre sum & \lrefx{Erdos269/ThreePrimeRunningLcm.lean}{384}{smoothHeightFiber_kernel_sum}\\
height-fibre normal form & \lrefx{Erdos269/ThreePrimeRunningLcm.lean}{406}{finiteSmoothKernelSum_groupedByHeight}\\
cubic majorant & \lrefx{Erdos269/ThreePrimeRunningLcm.lean}{430}{threePrimeHeight_le_cube}\\
origin value & \lrefx{Erdos269/ThreePrimeRunningLcm.lean}{443}{kernel_235_origin}\\
value at two & \lrefx{Erdos269/ThreePrimeRunningLcm.lean}{450}{kernel_235_two}\\
value at three & \lrefx{Erdos269/ThreePrimeRunningLcm.lean}{458}{kernel_235_three}\\
value at six & \lrefx{Erdos269/ThreePrimeRunningLcm.lean}{467}{kernel_235_six}\\
non-separation witness & \lrefx{Erdos269/ThreePrimeRunningLcm.lean}{479}{kernel_235_not_rankOne}\\
variable-base tail step & \lrefx{Erdos269/ThreePrimeRunningLcm.lean}{487}{tailStateStep}\\
expanded tail step & \lrefx{Erdos269/ThreePrimeRunningLcm.lean}{490}{tailStateStep_eq}\\
smooth exponent shell & \lrefx{Erdos269/ThreePrimeRunningLcm.lean}{500}{smoothExponentShell}\\
short-interval uniqueness & \lrefx{Erdos269/ThreePrimeRunningLcm.lean}{509}{exponent_unique_in_short_interval}\\
first-coordinate projection & \lrefx{Erdos269/ThreePrimeRunningLcm.lean}{543}{smoothExponentShell_card_le_dropFirst}\\
third-coordinate projection & \lrefx{Erdos269/ThreePrimeRunningLcm.lean}{585}{smoothExponentShell_card_le_dropThird}\\
sorted quadratic estimate & \lrefx{Erdos269/ThreePrimeRunningLcm.lean}{629}{sorted_pair_quadratic}\\
quadratic shell bound & \lrefx{Erdos269/ThreePrimeRunningLcm.lean}{646}{smoothExponentShell_card_quadratic}\\
dyadic internal power & \lrefx{Erdos269/ThreePrimeRunningLcm.lean}{662}{DyadicInternalPower}\\
internal-power uniqueness & \lrefx{Erdos269/ThreePrimeRunningLcm.lean}{668}{dyadicInternalPower_exponent_unique}\\
dyadic block base & \lrefx{Erdos269/ThreePrimeRunningLcm.lean}{688}{dyadicBlockBase235}\\
exact dyadic radix alphabet & \lrefx{Erdos269/ThreePrimeRunningLcm.lean}{699}{dyadicBlockBase235_cases}\\
bounded-radix consequence & \lrefx{Erdos269/ThreePrimeRunningLcm.lean}{711}{dyadicBlockBase235_mem_interval}\\
least positive residue & \lrefx{Erdos269/ResidueEscape.lean}{19}{leastPositiveResidue}\\
positive representative range & \lrefx{Erdos269/ResidueEscape.lean}{24}{leastPositiveResidue_pos_le}\\
representative congruence & \lrefx{Erdos269/ResidueEscape.lean}{45}{leastPositiveResidue_modEq}\\
escape condition & \lrefx{Erdos269/ResidueEscape.lean}{65}{ResidueEscapesWindow}\\
finite natural-state contradiction & \lrefx{Erdos269/ResidueEscape.lean}{71}{no_bounded_positive_state_of_residue_escape}\\
contrapositive form & \lrefx{Erdos269/ResidueEscape.lean}{92}{residue_le_bound_of_bounded_positive_state}\\
integer least-positive-residue consumer & \lrefx{Erdos269/ResidueEscape.lean}{104}{no_bounded_positive_int_state_of_leastPositiveResidue}\\
window base & \lrefx{Erdos269/RestrictedFloorSum.lean}{16}{windowBase}\\
window forcing & \lrefx{Erdos269/RestrictedFloorSum.lean}{21}{windowForcing}\\
affine window identity & \lrefx{Erdos269/RestrictedFloorSum.lean}{27}{affineRecurrence_window}\\
scaled forcing identity & \lrefx{Erdos269/RestrictedFloorSum.lean}{41}{windowForcing_const_mul}\\
integral carry window & \lrefx{Erdos269/RestrictedFloorSum.lean}{52}{integralCarry_window}\\
cofinal window producer & \lrefx{Erdos269/RestrictedFloorSum.lean}{69}{CofinalLocalWindowEscape}\\
reduced-carry extinction & \lrefx{Erdos269/RestrictedFloorSum.lean}{83}{no_positive_reducedCarry_of_cofinalLocalWindowEscape}\\
\end{tabular}
\fi

\begin{thebibliography}{9}
\small
\raggedright
\setlength{\itemsep}{0pt}
\bibitem{erdosgraham1980} P.~Erd\H{o}s and R.~L. Graham,
  \href{https://mathweb.ucsd.edu/~ronspubs/80_11_number_theory.pdf}{\emph{Old
  and New Problems and Results in Combinatorial Number Theory}},
  Monogr. Enseign. Math. 28, Geneva, 1980, p.~65.
\bibitem{erdos1988} P.~Erd\H{o}s, \emph{On the irrationality of certain series:
  problems and results}, in A.~Baker (ed.), \emph{New Advances in Transcendence
  Theory}, Cambridge UP, 1988, pp.~102--109,
  doi:\href{https://doi.org/10.1017/CBO9780511897184.009}{10.1017/CBO9780511897184.009}.
\bibitem{erdos1974letter} P.~Erd\H{o}s, letter to the editor (written 1~January
  1973), Fibonacci Quart. \textbf{12} (1974), p.~335.
\bibitem{erdosproblems} T.~F. Bloom,
  \href{https://www.erdosproblems.com/269}{\emph{Erd\H{o}s Problem \#269}},
  \texttt{erdosproblems.com/269}, accessed 22 July 2026.
\bibitem{formalconjectures269} The Formal Conjectures Authors,
  \href{https://github.com/google-deepmind/formal-conjectures/blob/main/FormalConjectures/ErdosProblems/269.lean}
  {\emph{FormalConjectures.ErdosProblems.\texttt{269}}},
  Lean source, 2025, accessed 27 July 2026.
\end{thebibliography}

\end{document}
